\section{训练与采样算法}

\subsection{训练}

式~\eqref{eq:l-simple} 已经把所有东西都准备好了，训练过程也就没有任何魔法：

\begin{algorithm}[htbp]
  \caption{扩散模型的训练（对应式~\eqref{eq:l-simple}）}
  \label{alg:train}
  \begin{algorithmic}[1]
    \Require 训练集 $\{\bx_0^{(n)}\}\sim p_{\text{data}}$，噪声调度 $\{\beta_t\}_{t=1}^{T}$，网络 $\beps_\theta$
    \Repeat
      \State 采样一个干净样本 $\bx_0\sim p_{\text{data}}$
      \State 采样一个噪声等级 $t\sim\mathcal{U}\{1,\dots,T\}$
      \State 采样一撮噪声 $\beps\sim\Ncal(\bzero,\bI)$
      \State 由闭式解合成含噪样本
             $\bx_t=\sqrt{\abt}\,\bx_0+\sqrt{1-\abt}\,\beps$
      \State 取梯度下降步，梯度为
             $\nabla_\theta\bigl\|\beps-\beps_\theta(\bx_t,t)\bigr\|^2$
    \Until{收敛}
  \end{algorithmic}
\end{algorithm}

注意算法~\ref{alg:train} 的一个特点：\emph{每一轮只加一次噪}。
因为有了闭式解，$t$ 可以任意跳，不需要按顺序走完整条链，
所以每一步的计算代价与 $T$ 无关。这是扩散模型训练成本其实并不高的原因。

\subsection{采样}

采样则是反向走完 $T$ 步：

\begin{algorithm}[htbp]
  \caption{扩散模型的采样（DDPM，$T$ 步）}
  \label{alg:sample}
  \begin{algorithmic}[1]
    \Require 训练好的网络 $\beps_\theta$，噪声调度 $\{\beta_t\}$，反向方差 $\sigma_t^2$
    \Ensure 一个生成样本 $\bx_0$
    \State 采样起点 $\bx_T\sim\Ncal(\bzero,\bI)$
    \For{$t=T,T-1,\dots,1$}
      \If{$t>1$}
        \State 采样 $\bz\sim\Ncal(\bzero,\bI)$
      \Else
        \State $\bz\gets\bzero$
      \EndIf
      \State $\bx_{t-1}\gets\dfrac{1}{\sqrt{\alpha_t}}
             \Bigl(\bx_t-\dfrac{\beta_t}{\sqrt{1-\abt}}\,\beps_\theta(\bx_t,t)\Bigr)
             +\sigma_t\,\bz$
    \EndFor
    \State \textbf{return} $\bx_0$
  \end{algorithmic}
\end{algorithm}

算法~\ref{alg:sample} 的结构与算法~\ref{alg:train} 高度对称：训练是“把已知的噪声加进去”再让网络
把它猜出来，采样是“把预测出来的噪声减掉”。把两段伪代码并排看一遍，
基本就理解扩散模型的全部机制了。

\subsection{反向方差 $\sigma_t^2$ 的取法：随机采样与确定性采样}

DDPM 论文给出两种取法\cite{ho2020ddpm}：
\begin{equation}
  \sigma_t^2=\beta_t
  \qquad\text{或}\qquad
  \sigma_t^2=\tilde\beta_t=\frac{1-\bar\alpha_{t-1}}{1-\abt}\,\beta_t .
  \label{eq:sigma-choices}
\end{equation}
前者在 $\abt$ 变化平缓时与后者接近；后者在 $t=1$ 时恰好为 $0$，意味着最后一步不加噪声。
$\sigma_t$ 实际上是一个“随机性旋钮”：取正数，采样是随机的；取 $\sigma_t=0$，
反向链就退化成一条\emph{确定性}的映射，同一个 $\bx_T$ 永远产生同一个输出。

\begin{remark}[DDIM：把 $1000$ 步压到 $50$ 步]
DDIM 从另一个角度出发\cite{song2021ddim}：它构造了一个\emph{非马尔可夫}的前向过程，
其边缘分布 $q(\bx_t\mid\bx_0)$ 与 DDPM 完全相同，因此可以复用同一个训练好的网络，
但采样式改写为“先估计 $\bx_0$、再重新构造 $\bx_{t-1}$”：
\begin{equation}
  \hat{\bx}_0=\frac{\bx_t-\sqrt{1-\abt}\,\beps_\theta(\bx_t,t)}{\sqrt{\abt}},
  \qquad
  \bx_{t-1}=\sqrt{\bar\alpha_{t-1}}\,\hat{\bx}_0+\sqrt{1-\bar\alpha_{t-1}}\,\beps_\theta(\bx_t,t).
  \label{eq:ddim}
\end{equation}
因为前向过程不再要求马尔可夫性，$\{\bx_t\}$ 可以只在任意\emph{子序列}
$\tau_1<\cdots<\tau_S$ 上取值。于是 $1000$ 步可以被压缩到 $20\sim50$ 步，
且\emph{同一个模型可以随意改变采样步数而无需重新训练}。
式~\eqref{eq:ddim} 与“取 $\sigma_t=0$”在形态上很接近（都没有再注入噪声），
但 DDIM 的跳步能力来自它的非马尔可夫结构，这一点是单纯的 $\sigma_t=0$ 做不到的。
\end{remark}

\subsection{为什么每一步要采样，而不是直接取均值？}

既然网络已经给出了均值 $\bmu_\theta$，每一步\emph{直接取均值}不是更省事吗？这是学习扩散模型时
很自然的一个疑问，它的答案有两层。

\textbf{第一层是概率上的。} $p_\theta(\bx_{t-1}\mid\bx_t)=\Ncal(\bmu_\theta,\sigma_t^2\bI)$
本身是一个\emph{分布}，均值只是它的中心。要“从分布中生成样本”，就必须从分布里采样；
若每一步都取均值，整条反向链变成一个确定性函数 $\bx_0=f(\bx_T)$，
同一个起点永远给出同一张图，众数唯一，样本多样性随之丧失。

\textbf{第二层是生成任务的实践经验。} 在序列生成（文本、语音）中，“每一步都取概率最大的那个词”
会得到重复、单调甚至崩塌的结果，这一现象早在语言模型采样中就有系统讨论
（\emph{The Curious Case of Neural Text Degeneration}\cite{holtzman2020curious}）；
语音合成中也观察到类似现象：完全去掉采样（例如去掉 dropout 带来的随机性）
会让韵律变得机械、不自然\cite{lee-ddpm}。

反过来说，这两点也说明\textbf{均值已经携带了绝大部分信息}。把 $\sigma_t$ 设成 $0$，
生成的图像依然可以接受、甚至更锐利，说明噪声更多影响的是多样性与细节的随机性，
而不是图像的主体结构。所以 $\sigma_t$ 更像是一个“多样性 $\leftrightarrow$ 确定性”的旋钮，
而不是一个非此即彼的选择。

\subsection{一个最小可运行的实现}

把上面两段伪代码直译成 PyTorch，核心不到 40 行（完整可运行版本见 \texttt{code/ddpm\_minimal.py}）：

\begin{lstlisting}[caption={DDPM 的最小实现（训练与采样）}, label={lst:ddpm}]
import torch

T = 1000
beta  = torch.linspace(1e-4, 2e-2, T)      # linear noise schedule
alpha = 1.0 - beta
abar  = torch.cumprod(alpha, dim=0)        # cumulative product of alpha

def q_sample(x0, t, noise):
    """Forward closed form: x_t = sqrt(abar_t) x0 + sqrt(1-abar_t) noise"""
    a = abar[t].sqrt().view(-1, *[1] * (x0.dim() - 1))
    s = (1.0 - abar[t]).sqrt().view(-1, *[1] * (x0.dim() - 1))
    return a * x0 + s * noise

def train_step(model, x0, opt):
    t     = torch.randint(0, T, (x0.shape[0],), device=x0.device)
    noise = torch.randn_like(x0)
    xt    = q_sample(x0, t, noise)              # one shot, O(1)
    loss  = ((model(xt, t) - noise) ** 2).mean() # predict the noise
    opt.zero_grad(); loss.backward(); opt.step()
    return loss.item()

@torch.no_grad()
def sample(model, shape):
    x = torch.randn(shape)                       # x_T ~ N(0, I)
    for ti in reversed(range(T)):
        t     = torch.full((shape[0],), ti, dtype=torch.long)
        eps   = model(x, t)
        mean  = (x - beta[ti] / (1.0 - abar[ti]).sqrt() * eps) / alpha[ti].sqrt()
        if ti > 0:
            x = mean + beta[ti].sqrt() * torch.randn_like(x)   # sigma_t^2 = beta_t
        else:
            x = mean
    return x
\end{lstlisting}
